\section*{Capítulo \ref{ch:b2:polymer-synthesis} --- Polímeros: síntesis, estructura y propiedades}

\begin{solution}{exo:b2:polymer-synthesis:1}
El polietileno tiene la unidad \ce{-[CH2-CH2]-}; el polipropileno, la unidad \ce{-[CH2-CH(CH3)]-}; el poli(cloruro de vinilo),
la unidad \ce{-[CH2-CHCl]-}; el poliestireno, la unidad \ce{-[CH2-CH(C6H5)]-}; y el nailon-6,6, la unidad \ce{-[NH(CH2)6NH-CO(CH2)4CO]-}.
\end{solution}

\begin{solution}{exo:b2:polymer-synthesis:2}
PET: por etapas (poliéster, con liberación de agua). Poliestireno: en cadena. Nailon-6,6: por etapas.
Poli(metacrilato de metilo): en cadena (un \omterm{def:b2:polymer-synthesis:polymer}{monómero} con \ce{C=C}). Poli(ácido láctico) a partir del hidroxiácido:
por etapas, un \omterm{def:b2:polymer-synthesis:polymer}{monómero} \ce{A-B}.
\end{solution}

\begin{solution}{exo:b2:polymer-synthesis:3}
$\bar M_n = 1000 \times \qty{104.15}{g/mol} \approx \qty{104}{kg/mol}$ (estireno \ce{C8H8},
\qty{104.152}{g/mol}).
% ledger: aw:C, aw:H
\end{solution}

\begin{solution}{exo:b2:polymer-synthesis:4}
\omterm{def:b2:polymer-synthesis:tacticity}{Sindiotáctica}. Sí: una cadena regular puede empaquetarse en cristalitos.
\end{solution}

\begin{solution}{exo:b2:polymer-synthesis:5}
Fracciones másicas 0.5 y 0.5. $1/\bar M_n = 0.5/10 + 0.5/100 = 0.055$, $\bar M_n \approx
\qty{18.2}{kg/mol}$; $\bar M_w = 0.5 \times 10 + 0.5 \times 100 = \qty{55}{kg/mol}$; $\text{\DH} \approx
3.0$.
\end{solution}

\begin{solution}{exo:b2:polymer-synthesis:6}
$1/(1-0.98) = 50$; $1/(1-0.995) = 200$.
\end{solution}

\begin{solution}{exo:b2:polymer-synthesis:7}
$r = 1/1.01$; para $p = 1$, $\bar X_n = (1+r)/(1-r) = (1.01 + 1)/(1.01 - 1) = 201$.
\end{solution}

\begin{solution}{exo:b2:polymer-synthesis:8}
$R_p \propto [\mathrm I]^{1/2}$: se multiplica por $\sqrt2 \approx 1.41$. $\nu \propto [\mathrm I]^{-1/2}$:
se divide por $\sqrt 2$, unas 0.71 veces su valor anterior.
\end{solution}

\begin{solution}{exo:b2:polymer-synthesis:9}
El PVC puro tiene su transición vítrea por encima de la temperatura ambiente: es un vidrio rígido. Las pequeñas moléculas de
plastificante se sitúan entre las cadenas, las separan y permiten que los segmentos se muevan a menor temperatura: $T_g$ baja
por debajo de la temperatura ambiente, y el material es gomoso y flexible en uso.
\end{solution}

\begin{solution}{exo:b2:polymer-synthesis:10}
$\dfrac{\mathrm d}{\mathrm dx}\bigl(x p^{\,x-1}\bigr) = p^{\,x-1}(1 + x\ln p)$, nula para $x = -1/\ln p$
(un máximo: el signo pasa de $+$ a $-$). Para $p$ cercano a 1, $\ln p \approx -(1-p)$, luego $x \approx
1/(1-p) = \bar X_n$.
\end{solution}

\begin{solution}{exo:b2:polymer-synthesis:11}
Etapa 1, \omterm{def:b2:acyl-substitution:transesterification}{transesterificación}: \ce{C6H4(COOCH3)2 + 2HOCH2CH2OH -> C6H4(COOCH2CH2OH)2 + 2CH3OH}, con el metanol
eliminado por destilación. Etapa 2: cada molécula de diéster lleva dos extremos \ce{CH2CH2OH}; un extremo ataca un éster
de otra molécula y libera una molécula de etano-1,2-diol, que se elimina por bombeo. El \omterm{def:b2:polymer-synthesis:polymer}{monómero}
lleva los dos tipos de grupo en la proporción correcta, como un \omterm{def:b2:polymer-synthesis:polymer}{monómero} \ce{A-B}: el equilibrio $r = 1$ está
incorporado, y eliminar el diol lleva $p$ hacia 1.
\end{solution}

\begin{solution}{exo:b2:polymer-synthesis:12}
Alimentación equimolar: razón $= (2.0 + 1)/(1 + 0.5) = 2$: el doble de unidades del \omterm{def:b2:polymer-synthesis:polymer}{monómero} 1. Con $r_1 = r_2 =
0$, razón $= [\mathrm M_1][\mathrm M_2]/([\mathrm M_2][\mathrm M_1]) = 1$ sea cual sea la alimentación: cada
radical adiciona solo el otro \omterm{def:b2:polymer-synthesis:polymer}{monómero}, un \omterm{def:b2:polymer-synthesis:copolymer}{copolímero} alternante.
\end{solution}

\begin{solution}{pb:b2:polymer-synthesis:1}
\textbf{1.} \ce{H2N(CH2)6NH2} y \ce{HOOC(CH2)4COOH}.
\textbf{2.} Una molécula de agua por enlace:
\[\mathrm{{-}COOH + H_2N{-} \longrightarrow {-}CO{-}NH{-} + H_2O}.\]
\textbf{3.} La sal contiene exactamente una diamina por cada diácido; pesar por separado dos líquidos o sólidos
nunca alcanzaría la exactitud que exige la ecuación de Carothers.
\textbf{4.} La unidad \ce{-[NH(CH2)6NH-CO(CH2)4CO]-}, de fórmula \ce{C12H22N2O2}, \qty{226.32}{g/mol}.
\textbf{5.} La \omterm{def:b2:polymer-synthesis:polymer}{unidad de repetición} contiene dos unidades procedentes de \omterm{def:b2:polymer-synthesis:polymer}{monómeros}: $M_0 = 226.32/2 = \qty{113.16}{g/mol}$.
\textbf{6.} Los carbonos de la diamina (6) y del diácido (6).
\textbf{7.} $\bar X_n = 50$, $\bar M_n = 50 \times 113.16 \approx \qty{5.66}{kg/mol}$.
\textbf{8.} $\bar X_n = 100$, $\bar M_n \approx \qty{11.3}{kg/mol}$.
\textbf{9.} Con el 99~\% de los grupos reaccionados, la cadena media solo tiene 100 unidades, demasiado corta para una
fibra resistente; cada paso adicional hacia la \omterm{def:b2:continuous-reactors:conversion}{conversión} completa duplica o triplica la longitud de cadena.
\textbf{10.} $\bar X_n = 10\,000/113.16 = 88.4$; $p = 1 - 1/88.4 \approx 0.989$.
\textbf{11.} Eliminando el agua: el fundido se calienta muy por encima de su punto de fusión, al final a presión
reducida, de modo que el equilibrio sigue desplazándose hacia la amida.
\textbf{12.} Una impureza monofuncional bloquea extremos de cadena, y cualquier impureza altera el equilibrio 1\,:\,1.
\textbf{13.} $r = 1/1.01 \approx 0.990$.
\textbf{14.} $\bar X_n = (1+r)/(1-r) = 201$; $\bar M_n \approx 201 \times 113.16 \approx
\qty{22.7}{kg/mol}$.
\textbf{15.} $\bar X_n = 1.9901/(1.9901 - 2 \times 0.9901 \times 0.99) \approx 1.9901/0.0297 \approx
67$.
\textbf{16.} Grupos amina: cuando se agotan los grupos ácido, todas las cadenas terminan en \ce{NH2}.
\textbf{17.} Convierte un extremo amina en una amida que ya no puede reaccionar: bloquea cadenas, lo que limita
$\bar M_n$ e impide que sigan creciendo cuando el \omterm{def:b2:polymer-synthesis:polymer}{polímero} vuelve a fundirse.
\textbf{18.} Para mantener una viscosidad del fundido reproducible en el hilado y el moldeo, que de otro modo
derivaría a medida que las cadenas siguen reaccionando en el fundido.
\textbf{19.} $\text{\DH} = 1 + p = 1.99$.
\textbf{20.} $\bar M_w = 1.99 \times 11.3 \approx \qty{22.5}{kg/mol}$.
\textbf{21.} Fracción en número $n_1 = 1 - p = 0.01$, el 1~\% de las moléculas; fracción másica $w_1 = (1-p)^2 =
10^{-4}$, el 0.01~\% de la masa.
\textbf{22.} Cerca de $x = -1/\ln 0.99 \approx 99.5$, es decir, cerca de $\bar X_n = 100$ (ejercicio 10).
\textbf{23.} $1000/226.32 = \qty{4.42}{mol}$ de \omterm{def:b2:polymer-synthesis:polymer}{unidades de repetición}, con dos enlaces amida cada una: \qty{8.84}{mol} de
agua, $8.84 \times 18.015 \approx \qty{159}{g}$.
\textbf{24.} Cadenas demasiado cortas dan una fibra débil; cadenas demasiado largas dan un fundido demasiado viscoso para pasar
por los finos orificios de la hilera.
\textbf{25.} $\bar X_n = 20\,000/113.16 = 176.7$ y $p = 1 - 1/176.7 = \boldsymbol{\approx 0.994}$.
\end{solution}
